% id: 26-12
% book: 베이직쎈 미적분1 · 01 함수의 극한
% section: 유형 012 함수의 극한의 도형에의 활용
% figure: tikz:fig-26-12
% answer: 
% answer_source: 
% variant_level: 0 (원본 전사)
% 이 파일은 build.mjs 가 items.json 에서 만든다. 고칠 때는 items.json 을 고친다.
\smash{\raisebox{1.5mm}{\dmkichul{베이직쎈 미적분1 26쪽 12번}}}
\noindent\begin{minipage}[t]{0.58\linewidth}\vspace{0pt}\raggedright 오른쪽 그림과 같이 함수 $y=x^2$의 그래프 위의 원점이 아닌 한 점~$\pt{P}$와 원점~$\pt{O}$를 지나며 중심이 $y$축 위에 있는 원 $C$가 있다. 점~$\pt{P}$가 함수 $y=x^2$의 그래프를 따라 원점~$\pt{O}$에 한없이 가까워질 때, 원의 중심~$\pt{C}$는 점 $(0{,}\mkern3mu\mathopen{}a)$에 한없이 가까워진다. 이때 $a$의 값을 구하시오.\end{minipage}\hfill\begin{minipage}[t]{0.4\linewidth}\vspace{0pt}\centering \resizebox{\ifdim\width>\linewidth\linewidth\else\width\fi}{!}{% 26-12(26쪽) — 곡선 y=x^2 위의 원점이 아닌 점 P 와 원점 O 를 지나며 중심 C 가 y축 위에 있는 원. 스캔 화질이 낮아 TikZ 로 다시 그림.
% 원문에 있는 것만: 라벨 y=x^2, P, C, O, x, y 와 중심의 점. 눈금·좌표값은 없다(답 a 가 그림에서 읽히지 않게).
\begin{tikzpicture}[x=0.36cm, y=0.36cm, line width=0.5pt]
  \draw[->, >=stealth, line width=0.7pt] (-4.3,0) -- (4.5,0) node[below=1pt] {$x$};
  \draw[->, >=stealth, line width=0.7pt] (0,-2.3) -- (0,7.9) node[left=1pt] {$y$};
  \node[below left=-2pt and -3pt, font=\small] at (0,0) {$\pt{O}$};
  % y=x^2
  \draw[line width=0.6pt, domain=-2.60:2.60, samples=110, smooth] plot (\x, {\x*\x});
  % 원점 O 와 점 P(t, t^2) 를 지나고 중심이 y축 위에 있는 원
  \draw[line width=0.6pt] (0,2.60125) circle (2.60125);
  \fill (0,2.60125) circle (2pt);
  \node[anchor=east, font=\small, inner sep=3pt] at (0,2.60125) {$\pt{C}$};
  \fill (2.05,4.2025) circle (2pt);
  \node[anchor=west, font=\small, inner sep=2pt] at (2.18,4.45) {$\pt{P}$};
  \node[anchor=west, font=\small] at (2.30,7.35) {$y=x^2$};
\end{tikzpicture}
}\end{minipage}\par
